\subsection{Converse and reduction}

\begin{theorem}[Aperiodic MLD classification --- the converse (Step 4)]\label{thm:aperiodic-mld-classification}\notready\uses{thm:mld-refines-cohequiv}
\emph{(Open frontier --- definitional residue.)} The substantive converse says
that cohomological equivalence implies geometric MLD.  The required carriers
below now exist: common-base classification
(\cref{thm:common-base-mld-classification}), cross-base up-to-closure
coinduction (\cref{thm:coalgebraic-crossbase-mld}), ladder climb
(\cref{thm:transfinite-ladder-climb}), chromatic non-idempotent carrier
(\cref{thm:chromatic-lorentzian-carrier}), aperiodic settling frame
(\cref{thm:aperiodic-settle-frame}), and geometric completeness
(\cref{thm:geometric-mld-complete}).  The remaining red fiber is purely
definitional: identify the frame's direct \(\mathrm{settle}\) readout with an
independent primitive transfinite \(\mathrm{closureOp}\) iteration at
\(\kappa\ge\omega_1\).
\end{theorem}

\begin{theorem}[MLD classification reduces to completeness]\label{thm:aperiodic-mld-classification-reduction}\lean{Infomaterial.Persistence.Aperiodic.aperiodic_mld_classification_reduction}\leanok\uses{thm:mld-refines-cohequiv}
For any local-derivability structure, \(\mathrm{Mld}\iff\mathrm{CohEquiv}\)
holds iff the structure is complete: equal \(\mathrm{cl\text{-}inv}\) already
yields \(\mathrm{Mld}\).  The node isolates the open mathematical content as
completeness of the geometric MLD carrier, not satisfiability of completeness in
some coarser relation.
\end{theorem}
